\section{Sintaks Dasar JavaScript}

JavaScript adalah bahasa pemrograman yang berjalan di browser dan (dengan Node.js) di server \cite{mdn-learn-js}. Sintaks dasar meliputi: variabel (\texttt{let}, \texttt{const}, \texttt{var}), tipe data (number, string, boolean, object, array), operator, kondisi (\texttt{if}, \texttt{switch}), dan perulangan (\texttt{for}, \texttt{while}). Fungsi didefinisikan dengan \texttt{function} atau arrow function \texttt{() => \{\}} \cite{js-info}.

Skrip JavaScript dapat disisipkan di halaman dengan tag \texttt{<script>} (inline atau \texttt{src="file.js"}). Penempatan skrip di akhir \texttt{<body>} atau penggunaan \texttt{defer} memastikan DOM sudah tersedia saat skrip dijalankan \cite{mdn-learn-js}.

\begin{lstlisting}[caption={Contoh JavaScript dasar}, basicstyle=\ttfamily\small, frame=single]
const nama = "Pemrograman Web";
let count = 0;
function tambah() { count++; return count; }
console.log(nama, tambah());
\end{lstlisting}
